% page 310 of the facsimile (image p-309 of the scan)
% block 154.9 x 250.6 mm ; left margin 8.0 mm
\pgc{-12.700mm}{0.000mm}{
\l{\cel{3}302}
\l{}
\l{\sou{kontradanso}.-\cel{2}Salono-danso, ek kin parti o figuri, exeku-}
\l{\cel{3}tata sucedante da du o quar grupi qui stacas inter-regarde}
\l{\cel{3}- DEFIS.}
\l{}
\l{\sou{kontrafaktar}. (\sou{trans}.) I. Imitar o +reproduktar kontre-lege}
\l{\cel{3}(la verko di ulu, e lua-detrimente). - II. Imitar fraude}
\l{\cel{3}objekto di qua la karaktero esas autentika. - FI.}
\l{}
\l{\sou{kontraforto}.- (\sou{arkitekt.)}\cel{2}Pilastro, saliajo masonura,}
\l{\cel{3}sur qua apogas su muro qua suportas kargajo. - EFIRS.}
\l{}
\l{\sou{kontraktar}. (\sou{trans.})\cel{2}I. (\sou{gram.})\cel{2}Reduktar (du vokali o}
\l{\cel{3}du silabi) ad un. - II. Reduktar a volumino min granda,}
\l{\cel{3}sen diminuteso di la maso : (l) (\sou{anat.}) per la kurtigo}
\l{\cel{3}di la fibri ; (2) (\sou{fiziko}) per la pluapudigo di la mo-}
\l{\cel{3}lekuli. - DEFIS.}
\l{}
\l{\sou{kontrakturo}.- (\sou{patol.})\cel{2}Stretiguro, kun rigidesko, qua even-}
\l{\cel{3}tas en muskulo pos konvulso, nevralgio, rumatismi, e c.DF}
\l{}
\l{\sou{kontralto}.- (\sou{muziko)}. Che la mulieri, la voco maxim basa.}
\l{\cel{3}- DEFIRS.}
\l{}
\l{\sou{kontrapunto}. (\sou{muziko)}. Formo di muziko klasika, en qua onu}
\l{\cel{3}asemblas, cirkum frazo precipua quan\cel{2}onu selektis kom}
\l{\cel{3}temo, parti acesora qui interunionas interseque en kom-}
\l{\cel{3}binuri hermoniala determinita, noto kontre noto. -DEFS.}
\l{}
\l{\sou{kontrastar.(netrans., kun}). (\sou{aludante du kozi)}. Interopoz-}
\l{\cel{3}esar tale ke lia interkompareso igas ica saliar. - DEFIRS.}
\l{}
\l{\sou{kontratar}.\cel{2}(\sou{netrans., pri}). Asumar, ad ulu (obligeso).-II.}
\l{\cel{3}Debatar, diskutar konvenciono solena inter du stati,de-}
\l{\cel{3}terminante, klauzope, to quo esos stipulita da singla de}
\l{\cel{3}la partopreninti.- III. (\sou{aludante plura personi}). Inter-}
\l{\cel{3}obligar per akto autentika. - DEFIRS.}
\l{}
\l{\sou{kon}travencar. (\sou{netrans., pri}). Agar kontre la preskripti}
\l{\cel{3}da regularo, da lego. - DEFIS.}
\l{}
\l{\sou{kontre}.\cel{2}Opoze ad (ulu, ulo).}
\l{}
\l{kontreskarpo.\cel{2}(tekn., militarto). Parto di la fosato di}
\l{\cel{3}fortifikuro qua orientizesas kontre la eskarpajo. -}
\l{}
\l{\sou{kontributar}. (\sou{trans., ad)}. - Kunlaborar, sua-parte, a verko}
\l{\cel{3}komuna. Adportar sua quoto a spenso komuna. - DEFIRS.}
\l{}
\l{\sou{kontricar}. (\sou{netrans., pri}). (\sou{religio}).Repentar pri la peko,}
\l{\cel{3}quan eventigas la chagreno ofensir Deo. - DEFIS.}
\l{}
\l{\sou{kontrolar}. (\sou{trans.}) I. Submisar a verifiko da la adminis-}
\l{\cel{3}trantaro.- II. Submisar ad exameno minucioza. - DEFIRS.}
\l{}
\l{kontroversar. (trans.) Diskutar kun ulu pri punto di dok-}
\l{\cel{3}trino. - DEFIS.}
}
